\chapter{The Consultant's Toolkit}

The blog has categories. WordPress requires it: you cannot publish a post without assigning it to at least one category. Many bloggers treat this as a chore, clicking the obvious option and moving on.

Cook has twenty-six categories. In the early years, they reveal a consultant's working life.

\bigskip

\section*{The Professional Landscape}

The category breakdown from 2008-2012 tells a story of someone solving problems across domains:

\begin{longtable}{p{3.00cm}p{3.00cm}p{3.00cm}}
\toprule
\textbf{Category} & \textbf{Posts (2008-2012)} & \textbf{\% of Era} \tabularnewline
\midrule
Uncategorized & 388 & 22\% \tabularnewline
Math & 357 & 21\% \tabularnewline
Software development & 327 & 19\% \tabularnewline
Statistics & 205 & 12\% \tabularnewline
Computing & 145 & 8\% \tabularnewline
Business & 126 & 7\% \tabularnewline
\bottomrule
\end{longtable}

The spread is remarkably even. Math, software development, and uncategorized posts each claim around 20\%. Statistics, computing, and business together account for another quarter. Not yet a math blog, but a professional's blog that includes math.

The category "Software development" contains posts like "Tips for learning regular expressions" and "Literate programming and statistics." The category "Business" contains posts like "Task switching" and "Organizational scar tissue."

These are the concerns of someone who works in organizations.

\bookfigure{ch02_category_treemap.png}{Category distribution}

\bigskip

\section*{The Tools of 2008}

The early archive preserves a snapshot of the consultant's toolkit.

\textbf{PowerShell} gets its own category (47 posts total). Cook wrote about PowerShell extensively from 2008-2012, with the heaviest concentration in the first two years. In February 2009, he linked to Joe Pruitt's "PowerShell A-Z" series. This was the tool of someone automating tasks in Windows environments. The last PowerShell post appeared in August 2012.

\textbf{Perl} appears in 20 posts across the blog's lifetime. "Periodic table of Perl operators" (2008). "Perl One-Liners Explained" (2012). "Learn one Perl command" (2012). Practical tools for practical problems. The last Perl post appeared in March 2021: "tcgrep: grep rewritten in Perl."

\textbf{Mathematica} gets 74 tag mentions across all years. It's a tool Cook has used consistently, though it never dominated. The early posts show someone who reached for whatever tool fit the problem.

The language distribution in the early era:

\begin{longtable}{l p{6.5cm}}
\toprule
\textbf{Language} & \textbf{Posts (2008-2012)} \tabularnewline
\midrule
Python & 77 \tabularnewline
PowerShell & 48 \tabularnewline
C++ & 33 \tabularnewline
Perl & 11 \tabularnewline
Mathematica & 9 \tabularnewline
\bottomrule
\end{longtable}

Python already leads in the early years, though it's far from dominant. By 2025, it would grow to 207 tagged posts.

\bigskip

\section*{The Clinical Trials Domain}

Fifty-five posts carry the "Clinical trials" category. They cluster in 2008-2013, then fade. By 2018, just one remained.

This was Cook's day job. From 2000 to 2013, he worked as a research statistician at MD Anderson Cancer Center, designing clinical trials for cancer treatments. He earned his PhD in mathematics from UT Austin in 1992, did a postdoc at Vanderbilt, and spent 1995-2000 in software development and project management before the MD Anderson role. The blog posts from this era reflect that work:

"Probability that a study result is true" (2008)
"Bayesian clinical trials in one zip code" (2009)
"Why drugs often list headache as a side-effect" (2010)

These posts were professional output---sharing methodology with colleagues, thinking through problems in public. They're technically dense. They assume an audience of statisticians working on similar problems.

When Cook left MD Anderson in 2013 to consult full-time, the clinical trials posts didn't follow. The topic was attached to a job, not to personal fascination.

The disappearance of clinical trials is evidence of what the blog would become: a place for what Cook chose to think about, not what he was paid to think about.

\bigskip

\section*{The Business of Consulting}

The early blog contains something the later blog lacks: business advice.

206 posts carry the "Business" category. Most appear before 2015. Titles include:

"How to avoid being outsourced or open sourced" (2008)
"Why programmers cannot be managed" (2008)
"Medieval software project management" (2008)
"Would you rather serve a market or a boss?" (2012)

These read like lessons learned. A consultant who has made mistakes and wants to save others the trouble. The voice is practical, not theoretical.

By the recent era, business posts nearly disappear. The transformation is complete: from consultant sharing professional wisdom to mathematician sharing mathematical fascination.

\bigskip

\section*{The Music Interlude}

Eighty-five posts carry the "Music" category. They span all seventeen years, but they never dominate.

Music posts reveal the blog's range:

"Circle of fifths and number theory" (2009): the connection between musical intervals and modular arithmetic.

"Typesetting music with LilyPond" (2009): the tools for creating sheet music programmatically.

"The acoustics of Hagia Sophia" (2018): the physics of sound in architectural spaces.

Not the posts of a professional musician. The posts of a mathematician who notices that music is also math. The same curiosity that drives the special functions posts: what is the structure underneath the surface?

The music category demonstrates that even in the early era, the blog followed fascination, not just profession. Music was never a job. It was always a choice.

\bigskip

\section*{The Geography of Curiosity}

Thirty posts carry the "Geography" category. Forty-two carry the "Orbital mechanics" tag.

These are the outliers: topics that appear, get explored, and persist at low frequency. They reveal a mind that finds mathematics everywhere:

"Finding distances using latitude and longitude" (2008)
"Mercator projection" (2009)
"Journey away from the center of the Earth" (2011)

The geography posts are typically short, practical, and include code. Someone needed to calculate something about the Earth, and Cook explained how.

Orbital mechanics posts similarly bridge abstract math and concrete application:

"Kepler and the contraction mapping theorem" (2018)
"Orbital resonance in Neptune's moons" (2019)
"Hohmann transfer orbit" (2020)

These posts demonstrate what the later blog would crystallize: mathematics as the universal tool, applicable to any domain that has structure.

\bigskip

\section*{What the Early Categories Reveal}

The 2008-2012 category distribution shows a generalist. Someone who works on many things, writes about all of them, and hasn't yet discovered what will persist.

Compare to the 2020-2025 distribution:

\begin{longtable}{p{3.00cm}p{3.00cm}p{3.00cm}}
\toprule
\textbf{Category} & \textbf{Posts (2020-2025)} & \textbf{\% of Era} \tabularnewline
\midrule
Math & 1,274 & 64\% \tabularnewline
Uncategorized & 305 & 15\% \tabularnewline
Computing & 245 & 12\% \tabularnewline
Statistics & 96 & 5\% \tabularnewline
Python & 22 & 1\% \tabularnewline
\bottomrule
\end{longtable}

Math has nearly doubled its share. Business has collapsed. Clinical trials has vanished. The consultant's toolkit has been replaced by the mathematician's focus.

The early era's breadth wasn't a mistake. It was exploration. The later era's focus emerged from that exploration---from five years of writing about everything, and discovering what kept being interesting.

\bigskip

\section*{The Uncategorized Mystery}

Nine hundred and ninety-six posts carry the "Uncategorized" label.

Some of these are genuinely miscellaneous: posts that don't fit anywhere. But many could have been categorized and weren't. The Uncategorized bin grows over time, suggesting that as the blog matured, the category system became less important.

In the early era, categories helped organize a scattered practice. In the later era, the practice had its own coherence. The categories became formality.

